% GENERATED FILE. Edit canonical lessons, then run npm run generate:books.
% canonical-source-hash: fnv1a64:f697d330060e5d9e
% canonical-lessons: JA-C471-butai, JA-C471-engi, JA-C471-hakushu-suru, JA-C471-gasshou-suru, JA-C471-sakkyoku-suru

\chapter{On Stage}
\label{ch:ja-on-stage}
\hlchaptermodality{471}

\begin{chapteropening}
\textbf{By the end of this chapter:} \emph{I can talk about a show on stage.}

Complete the last lesson of chapter 471: I can talk about a show on stage.
\end{chapteropening}

\section[\texorpdfstring{\ja{ぶたい} (butai)}{ぶたい (butai)}]{\ja{ぶたい} (butai) — a stage}

\label{lesson:JA-C471-butai}



\begin{quote}
Before the new one: say the Japanese for an enka ballad, then the Japanese for whistling.
\end{quote}

\subsection*{You'll want to know: \ja{ぶたい}}
\textbf{\ja{ぶたい}} — \emph{butai} — \textquotedblleft{}a stage\textquotedblright{}.

\begin{quote}
\textbf{\ja{ぶたいで} \ja{うたいました。}} — \emph{butai de utaimashita} — \textquotedblleft{}I sang on stage.\textquotedblright{}
\end{quote}

\begin{cousinweb}
\textbf{\ja{ぶたい}} (\emph{butai}) is Sino-Japanese (kango): two Chinese characters read with their on'yomi, readings borrowed from Chinese. They are \emph{bu}, \textquotedblleft{}dance\textquotedblright{} (Middle Chinese \emph{mjuX}) and \emph{tai}, \textquotedblleft{}pedestal\textquotedblright{} (\emph{doj}) (KANJIDIC2; Middle Chinese in Baxter's transcription).
\end{cousinweb}

\begin{grammarlens}[title={What you've built}]
One more word for describing what you lived and felt.
\end{grammarlens}

\subsection*{Your turn}
\begin{itemize}
\raggedright
  \item \emph{Say it:} \emph{butai}
  \item \emph{Say it:} \emph{butai}, once more
  \item \emph{Say it:} \emph{butai de utaimashita} — I sang on stage.
  \item \emph{Recall:} say the Japanese for an enka ballad, then the Japanese for whistling, then say \emph{butai} again
\end{itemize}

\begin{tcolorbox}[breakable,colback=teal!4,colframe=teal!35!black,title={Before you move on}]
What does \ja{ぶたい} mean? (\textquotedblleft{}A stage\textquotedblright{}.) Say it once more.
\end{tcolorbox}
\section[\texorpdfstring{\ja{えんぎ} (engi)}{えんぎ (engi)}]{\ja{えんぎ} (engi) — acting; a performance}

\label{lesson:JA-C471-engi}



\begin{quote}
Before the new one: say the Japanese for whistling, then the Japanese for a stage.
\end{quote}

\subsection*{You'll want to know: \ja{えんぎ}}
\textbf{\ja{えんぎ}} — \emph{engi} — \textquotedblleft{}acting; a performance\textquotedblright{}.

\begin{quote}
\textbf{\ja{あの} \ja{はいゆうの} \ja{えんぎは} \ja{すごいです。}} — \emph{ano haiyū no engi wa sugoi desu} — \textquotedblleft{}That actor's acting is amazing.\textquotedblright{}
\end{quote}

\begin{cousinweb}
\textbf{\ja{えんぎ}} (\emph{engi}) is Sino-Japanese (kango): two Chinese characters read with their on'yomi, readings borrowed from Chinese. They are \emph{en}, \textquotedblleft{}performance\textquotedblright{} (Middle Chinese \emph{yenX}) and \emph{gi}, \textquotedblleft{}skill\textquotedblright{} (\emph{gjeX}) (KANJIDIC2; Middle Chinese in Baxter's transcription).
\end{cousinweb}

\begin{grammarlens}[title={What you've built}]
One more word for describing what you lived and felt.
\end{grammarlens}

\subsection*{Your turn}
\begin{itemize}
\raggedright
  \item \emph{Say it:} \emph{engi}
  \item \emph{Say it:} \emph{engi}, once more
  \item \emph{Say it:} \emph{ano haiyū no engi wa sugoi desu} — That actor's acting is amazing.
  \item \emph{Recall:} say the Japanese for whistling, then the Japanese for a stage, then say \emph{engi} again
\end{itemize}

\begin{tcolorbox}[breakable,colback=teal!4,colframe=teal!35!black,title={Before you move on}]
What does \ja{えんぎ} mean? (\textquotedblleft{}Acting; a performance\textquotedblright{}.) Say it once more.
\end{tcolorbox}
\section[\texorpdfstring{\ja{はくしゅする} (hakushu suru)}{はくしゅする (hakushu suru)}]{\ja{はくしゅする} (hakushu suru) — to applaud; to clap}

\label{lesson:JA-C471-hakushu-suru}



\begin{quote}
Before the new one: say the Japanese for a stage, then the Japanese for acting.
\end{quote}

\subsection*{You'll want to know: \ja{はくしゅする}}
\textbf{\ja{はくしゅする}} — \emph{hakushu suru} — \textquotedblleft{}to applaud; to clap\textquotedblright{}.

\begin{quote}
\textbf{\ja{えんそうの} \ja{あと、みんな} \ja{はくしゅしました。}} — \emph{ensō no ato, minna hakushu shimashita} — \textquotedblleft{}After the performance, everyone applauded.\textquotedblright{}
\end{quote}

Polite \textbf{\ja{はくしゅします}}, \emph{hakushushimasu}; a \ja{する} verb. Past (polite): \textbf{\ja{かんきゃくが} \ja{はくしゅしました。}}, \emph{kankyaku ga hakushu shimashita}.

\begin{cousinweb}
\textbf{\ja{はくしゅする}} (\emph{hakushu suru}) is Sino-Japanese (kango) \emph{hakushu} with the native verb \emph{suru}, \textquotedblleft{}to do\textquotedblright{}: two Chinese characters read with their on'yomi, readings borrowed from Chinese. They are \emph{haku}, \textquotedblleft{}clap\textquotedblright{} (Middle Chinese \emph{phaek}) and \emph{shu}, \textquotedblleft{}hand\textquotedblright{} (\emph{syuwX}) (KANJIDIC2; Middle Chinese in Baxter's transcription).
\end{cousinweb}

\begin{grammarlens}[title={What you've built}]
One more word for describing what you lived and felt.
\end{grammarlens}

\subsection*{Your turn}
\begin{itemize}
\raggedright
  \item \emph{Say it:} \emph{hakushu suru}
  \item \emph{Say it:} \emph{hakushu suru}, once more
  \item \emph{Say it:} \emph{ensō no ato, minna hakushu shimashita} — After the performance, everyone applauded.
  \item \emph{Recall:} say the Japanese for a stage, then the Japanese for acting, then say \emph{hakushu suru} again
\end{itemize}

\begin{tcolorbox}[breakable,colback=teal!4,colframe=teal!35!black,title={Before you move on}]
What does \ja{はくしゅする} mean? (\textquotedblleft{}To applaud; to clap\textquotedblright{}.) Say it once more.
\end{tcolorbox}
\section[\texorpdfstring{\ja{がっしょうする} (gasshō suru)}{がっしょうする (gasshō suru)}]{\ja{がっしょうする} (gasshō suru) — to sing in a choir}

\label{lesson:JA-C471-gasshou-suru}



\begin{quote}
Before the new one: say the Japanese for acting, then the Japanese for to applaud.
\end{quote}

\subsection*{You'll want to know: \ja{がっしょうする}}
\textbf{\ja{がっしょうする}} — \emph{gasshō suru} — \textquotedblleft{}to sing in a choir\textquotedblright{}.

\begin{quote}
\textbf{\ja{がっこうで} \ja{みんなで} \ja{がっしょうしました。}} — \emph{gakkō de minna de gasshō shimashita} — \textquotedblleft{}We all sang together as a choir at school.\textquotedblright{}
\end{quote}

Polite \textbf{\ja{がっしょうします}}, \emph{gasshōshimasu}; a \ja{する} verb. Past (polite): \textbf{\ja{こどもが} \ja{がっしょうしました。}}, \emph{kodomo ga gasshō shimashita}.

\begin{cousinweb}
\textbf{\ja{がっしょうする}} (\emph{gasshō suru}) is Sino-Japanese (kango) \emph{gasshō} with the native verb \emph{suru}, \textquotedblleft{}to do\textquotedblright{}: two Chinese characters read with their on'yomi, readings borrowed from Chinese. They are \emph{ga}, \textquotedblleft{}join\textquotedblright{} (Middle Chinese \emph{hop}) and \emph{shō}, \textquotedblleft{}chant\textquotedblright{} (\emph{tsyhangH}) (KANJIDIC2; Middle Chinese in Baxter's transcription).
\end{cousinweb}

\begin{grammarlens}[title={What you've built}]
One more word for describing what you lived and felt.
\end{grammarlens}

\subsection*{Your turn}
\begin{itemize}
\raggedright
  \item \emph{Say it:} \emph{gasshō suru}
  \item \emph{Say it:} \emph{gasshō suru}, once more
  \item \emph{Say it:} \emph{gakkō de minna de gasshō shimashita} — We all sang together as a choir at school.
  \item \emph{Recall:} say the Japanese for acting, then the Japanese for to applaud, then say \emph{gasshō suru} again
\end{itemize}

\begin{tcolorbox}[breakable,colback=teal!4,colframe=teal!35!black,title={Before you move on}]
What does \ja{がっしょうする} mean? (\textquotedblleft{}To sing in a choir\textquotedblright{}.) Say it once more.
\end{tcolorbox}
\section[\texorpdfstring{\ja{さっきょくする} (sakkyoku suru)}{さっきょくする (sakkyoku suru)}]{\ja{さっきょくする} (sakkyoku suru) — to compose music}

\label{lesson:JA-C471-sakkyoku-suru}



\begin{quote}
Before the new one: say the Japanese for to applaud, then the Japanese for to sing in a choir.
\end{quote}

\subsection*{You'll want to know: \ja{さっきょくする}}
\textbf{\ja{さっきょくする}} — \emph{sakkyoku suru} — \textquotedblleft{}to compose music\textquotedblright{}.

\begin{quote}
\textbf{\ja{かのじょは} \ja{この} \ja{うたを} \ja{さっきょくしました。}} — \emph{kanojo wa kono uta o sakkyoku shimashita} — \textquotedblleft{}She composed this song.\textquotedblright{}
\end{quote}

Polite \textbf{\ja{さっきょくします}}, \emph{sakkyokushimasu}; a \ja{する} verb. Past (polite): \textbf{\ja{せんせいが} \ja{きょくを} \ja{さっきょくしました。}}, \emph{sensei ga kyoku o sakkyoku shimashita}.

\begin{cousinweb}
\textbf{\ja{さっきょくする}} (\emph{sakkyoku suru}) is Sino-Japanese (kango) \emph{sakkyoku} with the native verb \emph{suru}, \textquotedblleft{}to do\textquotedblright{}: two Chinese characters read with their on'yomi, readings borrowed from Chinese. They are \emph{sa}, \textquotedblleft{}make\textquotedblright{} (Middle Chinese \emph{tsak}) and \emph{kyoku}, \textquotedblleft{}melody\textquotedblright{} (\emph{khjowk}) (KANJIDIC2; Middle Chinese in Baxter's transcription).
\end{cousinweb}

\begin{grammarlens}[title={What you've built}]
One more word for describing what you lived and felt.
\end{grammarlens}

\subsection*{Your turn}
\begin{itemize}
\raggedright
  \item \emph{Say it:} \emph{sakkyoku suru}
  \item \emph{Say it:} \emph{sakkyoku suru}, once more
  \item \emph{Say it:} \emph{kanojo wa kono uta o sakkyoku shimashita} — She composed this song.
  \item \emph{Recall:} say the Japanese for to applaud, then the Japanese for to sing in a choir, then say \emph{sakkyoku suru} again
\end{itemize}

\begin{tcolorbox}[breakable,colback=teal!4,colframe=teal!35!black,title={Before you move on}]
What does \ja{さっきょくする} mean? (\textquotedblleft{}To compose music\textquotedblright{}.) Say it once more.
\end{tcolorbox}
